\documentclass[12pt, a4paper]{article}

\usepackage[a4paper, margin=1in]{geometry}

\usepackage[utf8]{inputenc}
\usepackage[mathscr]{euscript}
\usepackage{amssymb,amsmath,amsthm,amsfonts}
\usepackage[shortlabels]{enumitem}
\usepackage{multicol,multirow}
\usepackage{lipsum}
\usepackage{balance}
\usepackage{calc}
\usepackage[colorlinks=true,citecolor=blue,linkcolor=blue]{hyperref}
\usepackage{import}
\usepackage{xifthen}
\usepackage{pdfpages}
\usepackage{transparent}
\usepackage{listings}

\newcommand{\incfig}[2][1.0]{
    \def\svgwidth{#1\columnwidth}
    \import{./figures/}{#2.pdf_tex}
}

\newcommand{\incimg}[2][1.0]{
  \includegraphics[width=#1\columnwidth]{./figures/#2}
}

\newlist{enumproof}{enumerate}{4}
\setlist[enumproof,1]{label=\arabic*., parsep=1em}
\setlist[enumproof,2]{label=\arabic{enumproofi}.\arabic*., parsep=1em}
\setlist[enumproof,3]{label=\arabic{enumproofi}.\arabic{enumproofii}.\arabic*., parsep=1em}
\setlist[enumproof,4]{label=\arabic{enumproofi}.\arabic{enumproofii}.\arabic{enumproofiii}.\arabic*., parsep=1em}

\renewcommand{\qedsymbol}{\ensuremath{\blacksquare}}

\lstdefinestyle{cpp}{
  language=C++,
  commentstyle=\color{gray},
  keywordstyle=\color{orange},
  stringstyle=\color{Green},
  basicstyle={\ttfamily\footnotesize},
  breakatwhitespace=true,
  breaklines=true,
  captionpos=b,
  keepspaces=true,
  numbers=none,
  tabsize=2,
  showstringspaces=false,
  morekeywords=[3]{bool, wchar_t, char16_t, char32_t, size_t, ptrdiff_t, int, long, short, char, float, double, void, const, volatile, static, extern, register, inline, constexpr, noexcept, struct, class, union, enum, typename, template, using, namespace, auto, decltype, mutable, explicit, friend, virtual, override, final},
  keywordstyle=[3]\color{orange},
  morekeywords=[4]{std, string, vector, map, unordered_map, set, unordered_set, pair, tuple, optional, variant, unique_ptr, shared_ptr, make_unique, make_shared, thread, mutex, lock_guard, unique_lock, scoped_lock, condition_variable, condition_variable_any, atomic, memory_order, seq_cst, relaxed, acquire, release, acq_rel, cout, cin, cerr, endl, move, forward, ref, is_lock_free, shared_mutex, memory_order_seq_cst, empty, defer_lock, adopt_lock, queue},
  keywordstyle=[4]\color{blue},
  morekeywords=[5]{join, detach, lock, unlock, wait, notify_one, notify_all, load, store, exchange, compare_exchange_weak, compare_exchange_strong, fetch_add, fetch_sub, atomic_thread_fence, joinable, wait_and_pop, front, pop, push},
  keywordstyle=[5]\color{blue}
}\lstset{style=cpp}

\input{letterfonts}

\newcommand{\mytitle}{CS3211 Tutorial 2}
\newcommand{\myauthor}{github/omgeta}
\newcommand{\mydate}{AY 25/26 Sem 2}

\begin{document}
\raggedright
\footnotesize
\begin{center}
{\normalsize{\textbf{\mytitle}}} \\
{\footnotesize{\mydate\hspace{2pt}\textemdash\hspace{2pt}\myauthor}}
\end{center}
\setlist{topsep=-1em, itemsep=-1em, parsep=2em}
%%%%%%%%%%%%%%%%%%%%%%%%%%%%%%%%%%%%%%%%%%%%%%%%%%%%%%
%                      Begin                         %
%%%%%%%%%%%%%%%%%%%%%%%%%%%%%%%%%%%%%%%%%%%%%%%%%%%%%%
\begin{enumerate}[Q\arabic*.]
  \item 
    \begin{enumerate}
      \item Naively, we might expect $a=1, b=1$ but we get undefined behaviour due to the data races on $a,b$

      \item Atomics are significantly faster due to needing only a single instruction for addition. 

      \item \lstinline|seq_cst| is the slowest (using the slower \lstinline|xchg| instruction), \lstinline|acquire/release| and \lstinline|relaxed| take similar time (using the simpler \lstinline|mov| instruction)
    \end{enumerate}

  \item 
    \begin{enumerate}
      \item 0 $\rightarrow$ (a) $\rightarrow$ (b) $\rightarrow$ \{(p), (q)\}\\
        0 $\rightarrow$ (a) $\rightarrow$ \{(p), (q)\} $\rightarrow$ (b)\\
        0 $\rightarrow$ \{(p), (q)\} $\rightarrow$ (a) $\rightarrow$ (b)\\
        0 $\rightarrow$ \{(p)\} $\rightarrow$ (a) $\rightarrow$ \{(q)\} $\rightarrow$ (b)\\
        0 $\rightarrow$ \{(p)\} $\rightarrow$ (a) $\rightarrow$ (b) $\rightarrow$ \{(q)\}\\
        0 $\rightarrow$ (a) $\rightarrow$ \{(p)\} $\rightarrow$ (b) $\rightarrow$ \{(q)\}

      \item Regardless of memory ordering, there is a consistent coherence order for each atomic variable across all threads.
    \end{enumerate}

  \item 
    \begin{enumerate}
      \item cout always prints $1$: (a) happens-before (q) via (a)$\rightarrow_{sb}$(b)$\rightarrow_{sw}$(p)$\rightarrow_{sb}$(q)\\
        No change with non-atomics; same argument holds. 

        Without the while loop, the cout can print $0$, as there is no guarantee for a synchronizes-with from (b)$\rightarrow$(p)

      \item First cout prints $1,3,5$: (b) synchronizes-with (p) so atleast (a) happens-before (q), and later values are also possible\\ 
        Second cout prints $3,5$: (d) synchronizes-with (r) so atleast (c) happens-before (r) and later value is also possible

        Data races are possible between (c), (e) on the first cout, and (e) on the second cout: the writes are not guaranteed to always happens-before the couts.

      \item First cout prints $1,3,5$: (b) synchronizes-with (t) so atleast (a) happens-before (u), and later values are also possible\\ 
        Second cout prints $1,3,5$: (q) synchronizes-with (v) so atleast (p) happens-before (w) and later values are also possible

        Data races are possible between (p), (r) on the first cout, (r) on the second cout, and (a), (p), (r) 
    \end{enumerate}

  \item 
    \begin{enumerate}
      \item No, (q) is either before or after (p) in the coherence order. In either case, only one can read 0 and write 1.

      \item 0 $\rightarrow$ (a) $\rightarrow$ (b) $\rightarrow$ (p) $\rightarrow$ (q)\\
        0 $\rightarrow$ (a) $\rightarrow$ (p) $\rightarrow$ (b) $\rightarrow$ (q)\\
        0 $\rightarrow$ (a) $\rightarrow$ (p) $\rightarrow$ (q) $\rightarrow$ (b)\\
        0 $\rightarrow$ (p) $\rightarrow$ (a) $\rightarrow$ (b) $\rightarrow$ (q)\\
        0 $\rightarrow$ (p) $\rightarrow$ (a) $\rightarrow$ (q) $\rightarrow$ (b)\\
        0 $\rightarrow$ (p) $\rightarrow$ (q) $\rightarrow$ (a) $\rightarrow$ (b)\\
        For all modification orders, the final value is $4$

      \item Yes, since we have two different modification orders per-atomic and no happens-before across threads contradicting the statement.\\
        x’s modification order: 0 (initialised) $\rightarrow$ 1 (q) $\rightarrow$ 2 (a)\\
        y’s modification order: 0 (initialised) $\rightarrow$ 1 (b) $\rightarrow$ 2 (p)

    \end{enumerate}
\end{enumerate}
%%%%%%%%%%%%%%%%%%%%%%%%%%%%%%%%%%%%%%%%%%%%%%%%%%%%%%
%                       End                          %
%%%%%%%%%%%%%%%%%%%%%%%%%%%%%%%%%%%%%%%%%%%%%%%%%%%%%%

\end{document}
